\chapter[\emj{🎨}~Grafo Bipartito (2-Coloring)]{\emj{🎨}~Grafo Bipartito (Coloreo con 2 Colores)}

\begin{dsaquees}

Repartir a un grupo de personas en DOS equipos 🔴🔵 de modo que cada
persona quede en equipo CONTRARIO al de todos sus rivales. Si logras
hacerlo sin conflictos, el grafo es bipartito.

El truco: pinta a alguien de rojo, obliga a todos sus vecinos a ser
azules, a los vecinos de esos a ser rojos, y así. Si en algún momento
alguien DEBE ser rojo y azul a la vez (un vecino ya tiene tu mismo
color), es imposible: el grafo NO es bipartito.

\end{dsaquees}

\begin{dsaotraforma}

Un grafo es \textbf{bipartito} si sus vértices se pueden partir en dos
conjuntos disjuntos tal que toda arista conecta un vértice de un
conjunto con uno del otro (ninguna arista queda dentro del mismo
conjunto).

Se verifica con un coloreo de 2 colores por BFS o DFS: asigna un color
al nodo inicial y el color OPUESTO a cada vecino. Si encuentras una
arista entre dos nodos del MISMO color, no es bipartito.

Teorema clave: un grafo es bipartito \textbf{si y solo si} no contiene
ningún ciclo de longitud impar.

\end{dsaotraforma}

\begin{dsadiagrama}
\begin{verbatim}
BIPARTITO (ciclo par)          NO BIPARTITO (ciclo impar)

  A(R) ─── B(A)                  A(R) ─── B(A)
   │        │                     │  ╲      │
  D(A) ─── C(R)                   │   ╲     │
                                 D(A) ─ C(R)? A-C exige A(R)-C debe ser
  2 equipos: {A,C} y {B,D}       azul, pero C ya conecta con B(A)...
  Toda arista cruza -> SÍ        triangulo A-B-C (ciclo 3) -> NO
\end{verbatim}
\end{dsadiagrama}

\begin{lstlisting}[language=C++]
CÓMO FUNCIONA (BFS coloring)
1. color[] = -1 (sin colorear) para todos.
2. Por cada componente no coloreada: colorea el inicio con 0, mételo a
   la Queue.
3. Al sacar u, cada vecino v: si no tiene color, dale el opuesto
   (1 - color[u]); si ya tiene el MISMO color que u -> NO bipartito.

📝 PSEUDOCÓDIGO
──────────────────────────────────────────────────────────────

función esBipartito(V, adj):
    color = arreglo de -1
    PARA cada nodo s sin color:
        color[s] = 0; q = Queue con s
        MIENTRAS q no vacia:
            u = q.dequeue()
            PARA cada vecino v de u:
                SI color[v] == -1:
                    color[v] = 1 - color[u]; q.enqueue(v)
                SINO SI color[v] == color[u]:
                    return FALSO
    return VERDADERO

💻 CÓDIGO C++
──────────────────────────────────────────────────────────────

#include <vector>
#include <queue>
using namespace std;

bool isBipartite(int V, vector<vector<int>>& adj) {
    vector<int> color(V, -1);
    for (int s = 0; s < V; s++) {
        if (color[s] != -1) continue;      // ya visitado
        color[s] = 0;
        queue<int> q; q.push(s);
        while (!q.empty()) {
            int u = q.front(); q.pop();
            for (int v : adj[u]) {
                if (color[v] == -1) {
                    color[v] = 1 - color[u];   // color opuesto
                    q.push(v);
                } else if (color[v] == color[u]) {
                    return false;              // conflicto
                }
            }
        }
    }
    return true;
}

⏱️ COMPLEJIDAD (Big-O)
──────────────────────────────────────────────────────────────

  Recurso     │ Costo
  ────────────┼─────────────
  Tiempo      │ O(V + E)
  Espacio     │ O(V)

* Recorre cada nodo y arista una vez; el array color es O(V).
\end{lstlisting}

\begin{dsaentrevistas}

✅ "Is Graph Bipartite?" (LeetCode 785) y "Possible Bipartition" (886):
   ambos son 2-coloring puro. Recuerda recorrer TODAS las componentes
   (el grafo puede estar desconectado).

💡 Aplicación real: asignar tareas a 2 máquinas, detectar conflictos de
   horarios, o modelar relaciones "esto vs aquello" (matching bipartito).

⚠️ Grafo desconectado: si olvidas el bucle externo sobre nodos sin
   color, solo verificas una componente y das un falso positivo.

🔑 Ciclo impar = no bipartito: si te preguntan la intuición, esa es la
   frase. Un triángulo (ciclo de 3) nunca es bipartito.

\end{dsaentrevistas}

\begin{dsaresumen}

• Bipartito: partir los nodos en 2 grupos; toda arista cruza entre ellos.
• Se verifica coloreando con 2 colores por BFS o DFS.
• Conflicto: una arista une dos nodos del mismo color -> no bipartito.
• Equivalente: no tiene ciclos de longitud impar.
• Recorre TODAS las componentes (grafo puede estar desconectado).
• O(V + E).
════════════════════════════════════════════════════════════════

\end{dsaresumen}
